\documentclass[12pt, a4paper]{report}
\usepackage[utf8]{inputenc}
\usepackage[T1]{fontenc}
\usepackage[english]{babel}
\usepackage{geometry}
\geometry{margin=2.5cm}

\begin{document}

\tableofcontents
\newpage

\begin{abstract}
% Résumé de 15 lignes maximum
\end{abstract}

\chapter*{Introduction}
\addcontentsline{toc}{chapter}{Introduction}
% Contexte général, motivation, objectifs du TER

\chapter{Background and Stakes}

\section{Why Model Pollution Spatially?}
% Interpolation, risk zones, public policy decisions

\section{Limitations of Non-Spatial Approaches}
% Why classical regression is insufficient

\section{Research Objectives}
% Pollutant choisi, zone géographique, questions de recherche

\chapter{Theoretical Framework: Spatial Statistics}

\section{Key Concepts}
\subsection{Stationarity and Isotropy}
\subsection{Spatial Autocorrelation}
% Moran's index, variogram

\section{Kriging}
\subsection{Ordinary Kriging}
\subsection{Universal Kriging}
\subsection{Kriging Variance and Uncertainty}

\section{Spatial Regression Models}
\subsection{Spatial Autoregressive Model (SAR)}
\subsection{Spatial Error Model (SEM)}
\subsection{Model Comparison}

\section{Further Extensions (optional)}
% SPDE/INLA, Bayesian approach

\chapter{Data}

\section{Sources}
% Copernicus, Atmo, OpenAQ, INSEE

\section{Study Area and Pollutant}
% Justification du choix

\section{Data Processing}
% Missing values, irregular sensor networks

\chapter{Numerical Application}

\section{Exploratory Spatial Analysis}
% Moran's I, variogram fitting

\section{Kriging Interpolation}
% Concentration maps, cross-validation

\section{Spatial Regression}
% SAR vs SEM, covariates

\section{Risk Assessment}
% Exceedance probability, spatial VaR

\section{Model Comparison and Validation}
% RMSE, cross-validation metrics

\chapter{Discussion}

\section{Main Results}
\section{Limits}
% Non-stationarity, edge effects, data quality

\section{Perspectives}
% Spatio-temporal models, link with health/actuarial risk

\chapter{Project Management}
% Concepts de gestion de projet appliqués au TER

\chapter*{Conclusion}
\addcontentsline{toc}{chapter}{Conclusion}

\bibliographystyle{plain}
\bibliography{references}

\appendix

\chapter{Group Organization and Interactions with Supervisor}
% Répartition des tâches, réunions, échanges avec l'encadrant

\chapter{AI Prompt Usage}
% Prompts utiles, prompts infructueux, démarcation IA / travail personnel

\listoffigures
\listoftables

\end{document}